\begin{lemma}[The PE hierarchy base strictly extends the terminal ladder tree]
Let $G$ be a finite loopless multigraph with distinguished edge $e_0$.
Let $\mathcal T$ be an extended Tashkinov tree whose generating coloring is
$\varphi$, let $n$ be its terminal rung, and put
\[
                         T_n=\mathcal T.\mathrm{tree}(n).
\]
Let $(T_{n,0}^{*},R_n,v_n)$ be a
\noderef{HierarchyBase} for $\mathcal T$ under the same coloring
$\varphi$.  If the terminal extension kind is PE, then
\[
 T_n\longrightarrow T_{n,0}^{*}
 \quad\text{is a $\varphi$-\noderef{TashkinovAugmentationSequence},}
 \qquad |T_n|<|T_{n,0}^{*}|.
\]
The lengths here count edge entries of the ordered tree-sequences.
\end{lemma}
\begin{proof}
Let $P=T_n(v_n)$ denote the exact prefix ending at the positive support
position stored in the hierarchy base.  The PE clause of
\noderef{HierarchyBase} gives a $\varphi$-Tashkinov closure from $P$ to
$R_n$ and an exact \noderef{OrderedTreeUnion}
\[
                         T_{n,0}^{*}=T_n\vee R_n.                \tag{1}
\]
By \noderef{TashkinovClosure}, this closure consists of an augmentation
sequence to $R_n$ together with closedness of $V(R_n)$.
The prefix $P$ is an exact ordered initial segment of $T_n$.  Apply
\noderef{OrderedTreeUnionAugmentationSequence} to the closure's
augmentation sequence $P\longrightarrow R_n$, this initial-segment
relation, and (1).  It gives
\[
                         T_n\longrightarrow T_{n,0}^{*}          \tag{2}
\]
as a $\varphi$-Tashkinov augmentation sequence.  It remains to prove
that (2) is nonempty.

The ETT has positive rung, so its terminal rung has a predecessor $q$.
Unpack the terminal certificate supplied by \noderef{TashkinovSeries}.
The displayed PE-kind hypothesis rules out the initial, SE, and RE
constructors of \noderef{TashkinovSeriesStep}.  We therefore use its PE
constructor.  Write
\(
 \pi
\)
for the selected maximum-defective coloring, $f_n$ for the selected exact
boundary-edge identity, $\delta_n$ for its color, $\gamma_n$ for the PE
extension color, $\rho$ for the second stable coloring, and
\[
 K=K_{\rho}(v_n;\gamma_n,\delta_n)
\]
for the \noderef{TwoColorComponent} switched by that constructor.  Thus the
terminal coloring is $\rho/K$.  Since the ETT's generating coloring is
$\varphi$, the generating-coloring equality in
\noderef{ExtendedTashkinovTree} identifies
\[
                            \varphi=\rho/K.                       \tag{3}
\]

We next construct a surviving $\delta_n$-boundary identity; its existence
is not an assumption.  The certificate \noderef{MaximumDefectiveChoice}
says that $f_n$ is a defective boundary identity under $\pi$ and that
\(\pi(f_n)=\delta_n\).  By \noderef{DefectiveEdge}, there is a defective
color $c$ for which $f_n$ lies in the exact $c$-boundary.  Membership in
that boundary, by \noderef{BoundaryColorEdges}, also says
\(\pi(f_n)=c\).  Hence $c=\delta_n$.  The definition
\noderef{DefectiveColor} now gives
\[
                    |\partial_{\pi,\delta_n}(T_n)|\geq2.         \tag{4}
\]
Choose an exact identity
\[
 g\in\partial_{\pi,\delta_n}(T_n),\qquad g\ne f_n.              \tag{5}
\]

Every boundary identity is incident with a tree vertex by
\noderef{BoundaryEdges} and hence is a
\noderef{TreeIncidentEdge}.  The second stable relation in the PE
constructor is a \noderef{StableColoring} whose control set contains
$\delta_n$.  Its edge clause therefore applies separately to the exact
identities $f_n$ and $g$, and (5) gives
\[
                         \rho(f_n)=\rho(g)=\delta_n.              \tag{6}
\]
The PE rule \noderef{ParallelExtensionColor}, together with the vertex
clause of the same stable relation, says that $\gamma_n$ is missing at
$v_n$ under $\rho$.  Since $f_n$ is incident with $v_n$ and has color
$\delta_n$ by (6), \noderef{MissingColors} and
\noderef{PresentColors} give $\gamma_n\ne\delta_n$.

The incidence rule in \noderef{TwoColorIncidenceGraph} and (6) put the
edge-node of $f_n$ in $K$.  Expanding
\noderef{TwoColorComponentSwap}, the switch changes this exact identity
from $\delta_n$ to $\gamma_n$:
\[
                         (\rho/K)(f_n)=\gamma_n.                  \tag{7}
\]
We claim that the edge-node of $g$ is not in $K$.  If it were, (6) and
the same switch formula would give $(\rho/K)(g)=\gamma_n$.  Boundary
membership is independent of coloring, so both $g$ and $f_n$ would then
belong to the terminal $\gamma_n$-boundary.  But
\noderef{HkuParallelBoundaryStructure} says that this exact boundary is
the singleton $\{f_n\}$.  This would force $g=f_n$, contrary to (5).
Thus $g$ is outside $K$, and the switch fixes it.  Using (3),
\[
                   g\in\partial_{\varphi,\delta_n}(T_n).         \tag{8}
\]

For completeness, $\delta_n$ is missing at $v_n$ after the switch.  If
an edge identity incident with $v_n$ had terminal color $\delta_n$, then
an identity inside $K$ would have had color $\gamma_n$ under $\rho$,
contradicting that $\gamma_n$ is missing at $v_n$.  An identity outside
$K$ would have had color $\delta_n$ under $\rho$, but its incidence with
$v_n$ would itself put its edge-node in $K$ by
\noderef{TwoColorIncidenceGraph}.  Both alternatives are impossible, so
\[
                       \delta_n\in\overline\varphi(v_n).         \tag{9}
\]

Apply the $\delta_n$-location clause of
\noderef{HkuParallelBoundaryStructure} to the now-established identity
(8).  There is a vertex $w\prec_{T_n}v_n$ incident with $g$.  By the
definition of the positive prefix in \noderef{TreeSequencePrefix}, both
$w$ and $v_n$ lie in $V(P)$.  The closure augmentation sequence from
$P$ to $R_n$ is an exact initial extension by
\noderef{TashkinovAugmentationSequenceInitialSegment}.  Consequently
\[
                              w,v_n\in V(R_n).                   \tag{10}
\]

Since $g$ is a boundary identity of $T_n$, write its two ends as
$w'$ and $z$, where $w'\in V(T_n)$ and $z\notin V(T_n)$.  The vertex
$w$ is incident with $g$ and belongs to $V(T_n)$, so $w=w'$; the other
possibility would put the outside end $z$ in $T_n$.  Suppose for a
contradiction that $z\notin V(R_n)$.  Then (10) makes $g$ a boundary
identity of $R_n$.  Equations (8) and (3) give $\varphi(g)=\delta_n$,
while (9) and (10) give a vertex $v_n\in V(R_n)$ missing
$\delta_n$.  This contradicts the closedness clause of the Tashkinov
closure in \noderef{HierarchyBase}, namely
\noderef{ColoredVertexSetClosed}.  Hence
\[
                        z\in V(R_n)-V(T_n).                      \tag{11}
\]

The vertex-set clause of the exact ordered union (1) is
\[
 V(T_{n,0}^{*})=V(T_n)\cup V(R_n).
\]
Thus (11) gives a vertex of $T_{n,0}^{*}$ which is not a vertex of
$T_n$.  On the other hand, (2) and
\noderef{TashkinovAugmentationSequenceInitialSegment} make $T_n$ an exact
ordered initial segment of $T_{n,0}^{*}$.  If their lengths were equal,
entrywise equality in \noderef{TreeSequenceInitialSegment}, together
with vertex injectivity from \noderef{TreeSequence}, would make their
vertex sets equal, contradicting (11).  Therefore
$|T_n|<|T_{n,0}^{*}|$, as required.
\end{proof}
